\documentclass{article}

% Standard Packages
\usepackage[utf8]{inputenc}
\usepackage[T1]{fontenc}
\usepackage{hyperref}
\usepackage{url}
\usepackage{booktabs}
\usepackage{amsfonts}
\usepackage{nicefrac}
\usepackage{microtype}
\usepackage{xcolor}
\usepackage{graphicx}
\usepackage{amsmath}
\usepackage{caption}
\usepackage{subcaption}
\usepackage{geometry}
\geometry{letterpaper, margin=1in}

\title{Spectral Dynamics of Attention Under Cognitive Load: \\ Analyzing the Low-Rank Dissipation of MHA Across Bloom's Taxonomy}

\author{
  Vedant Maheshwari \\
  \texttt{vedant.maheshwari@example.com} \\
}
\date{}

\begin{document}

\maketitle

\begin{abstract}
As large language models scale, understanding the intrinsic dimensionality and behavior of Multi-Head Attention (MHA) layers has become an increasingly urgent priority. While MHA is known to construct rich semantic representations, the structural necessity of its full dimensional space during inference is heavily debated. This paper presents an empirical investigation into the spectral properties of MHA under varying levels of cognitive complexity, systematically guided by Bloom’s Taxonomy of Learning (BTL). We analyze the Phi-2 language model using Singular Value Decomposition (SVD) on attention matrices to map the dissipation of energy, effective rank, and spectral entropy across multiple network layers. Our findings indicate that MHA functions as a strongly dissipative, low-rank system, which severely compresses representations in deeper layers. By truncating attention structures via SVD—maintaining only the top-$k$ singular components during causal generation—we reveal that downstream reasoning degradation correlates significantly with higher cognitive BTL tasks (e.g., Evaluating, Creating). These results provide a robust mathematical lens explaining why factual recall tasks persist under drastic rank ablation, while generative synthesis collapses, offering a pathway toward adaptive-rank inference optimization based on cognitive load.
\end{abstract}

\section{Introduction}

Multi-Head Attention (MHA) serves as the computational backbone of modern sequential models. Despite its ubiquitous adoption, calculating exact $Q K^T$ interactions natively scales quadratically with sequence length, enforcing massive memory constraints during inference. Several approaches, such as Grouped Query Attention (GQA), seek to compress this behavior, yet questions remain regarding the inherent information density encoded inside the standard self-attention matrices during distinct types of reasoning.

While researchers often analyze MHA via attention rollout maps or isolated head induction, we propose examining the structural energy of the attention probability matrix as a continuous linear projection. By treating the transition matrix as a linear operator and evaluating its Singular Value Decomposition (SVD), we can mathematically decouple the "dominant semantics" encoded by the primary singular vectors from the "multi-scale nuances" encapsulated in the long-tail singular vectors.

Crucially, real-world inference prompts heavily oscillate in complexity. Providing the exact same full-rank computational bandwidth to answer "What is the capital of France?" compared to "Design a novel algorithmic sorting technique" is likely sub-optimal. This paper utilizes Bloom’s Taxonomy of Learning (BTL)—a pedagogical framework categorizing cognitive load into six hierarchical stages—to structurally evaluate the resilience of attention layers.  

We summarize our contributions as follows:
\begin{enumerate}
    \item We introduce a targeted, multi-layer SVD intervention framework inside the causal generation loop of Phi-2, dynamically ablating trailing singular components across inference.
    \item We evaluate variations in Effective Rank ($r_{\text{eff}}$), Spectral Entropy ($H$), and Spectral Energy across the depth of the transformer, verifying that MHA acts as a \textit{low-rank dissipator}.
    \item We execute a comprehensive performance sweep across 42 BTL-aligned Chain-of-Thought (CoT) prompts. We demonstrate that fundamental recall tasks survive severe rank ablation ($k=1$), while synthetic and evaluative reasoning depend catastrophically on higher-rank representations.
\end{enumerate}

\section{Methodology and Experimental Setup}

\subsection{Singular Value Decomposition (SVD) Intervention Pipeline}

To systematically investigate the dependency on matrix dimensionality, we intercepted the standard attention forward pass of the Phi-2 model. Rather than allowing the softmax attention weights $A \in \mathbb{R}^{Q \times K}$ to guide value projection natively, we factorize the matrix and enforce a hard rank threshold. 

For each attention head, let $A$ be the post-softmax attention probability matrix. We apply SVD such that $A = U \Sigma V^T$, where $\Sigma$ is the diagonal matrix of singular values $\sigma_1 \ge \sigma_2 \ge \dots \ge \sigma_n > 0$. We construct a lower-rank approximation $A_k$ by keeping only the top-$k$ singular components:

\begin{equation}
    A_k = \sum_{i=1}^{k} \sigma_i u_i v_i^T = U_k \Sigma_k V_k^T
\end{equation}

Where $U_k$ and $V_k$ consist of the first $k$ columns of $U$ and $V$. We perform this ablation dynamically inside layers 15 through 20 (typically representing middle-level reasoning distribution) and layer 31 (the final layer). The truncated matrix $A_k$ is then multiplied against the value states $V_{states}$ to resume the autoregressive step.

\subsection{Quantifying Spectral Properties}

To monitor the structural density of information flowing through the intervened network, we continuously record multiple information-theoretic metrics for every batch calculation:

\paragraph{Energy Retention} 
The square of the singular values represents the spectral energy. The percentage of energy retained after a $k$-truncation is bounded by:
\begin{equation}
    E_{\text{ret}} = \frac{\sum_{i=1}^k \sigma_i^2}{\sum_{j=1}^n \sigma_j^2} \times 100
\end{equation}

\paragraph{Spectral Entropy and Effective Rank}
The singular values denote a discrete probability distribution over the matrix's structural eigen-directions. By normalizing the eigenvalues, we obtain a spectral density profile:
\begin{equation}
    p_i = \frac{\sigma_i^2}{\sum_{j=1}^n \sigma_j^2}
\end{equation}
The Shannon Information Entropy over this spectral distribution defines the matrix's diffusion of information, or \textit{Spectral Entropy} ($H$):
\begin{equation}
    H = - \sum_{i=1}^{n} p_i \ln(p_i)
\end{equation}
From this entropy, we define a continuous metric of dimensionality known as the \textit{Effective Rank} ($r_{\text{eff}}$), formally:
\begin{equation}
    r_{\text{eff}} = \exp(H)
\end{equation}

\subsection{Bloom's Taxonomy Level (BTL) Evaluation Suite}

We designed an evaluation suite consisting of 42 distinct reasoning prompts. These prompts are uniformly distributed across the six hierarchical boundaries of Bloom's Taxonomy: \textit{Remembering, Understanding, Applying, Analyzing, Evaluating,} and \textit{Creating}. All prompts are formatted as Instruct-based Chain-of-Thought directives (e.g., ``Think step-by-step to arrive at a conclusion.''). During validation, we measure the KL-divergence of the ensuing token logit distributions and verify generative identicality metrics compared to unmodified baseline outputs.

\section{Experimental Results and Multi-Layer Analysis}

\subsection{Effective Rank Diminishes with Network Depth}

We observed deep structural changes in representation moving from intermediate reasoning layers (L15-L20) to the output boundary (L31). As visualized in our findings:

\begin{itemize}
    \item \textbf{Middle Layers (15-20):} Demonstrated higher levels of structural variance. The Effective Rank averaged between $2.5$ and $4.0$. This higher dimensionality signals the active manipulation of multi-modal features, where attention heads must track multiple distinct interacting concepts simultaneously.
    \item \textbf{Final Layer (31):} Effective Rank collapses significantly to $r_{\text{eff}} \approx 1.5$, regardless of sequence length. At this depth, over $~95\%$ of spectral energy consolidates uniformly into the singular primary component ($\sigma_1$).
\end{itemize}

This supports the hypothesis that multi-head attention acts as a strongly dissipative bottleneck throughout processing. MHA ingests high-rank combinations of token interactions inside intermediate layers but sinks variance systematically prior to the un-embedding projection. 

\subsection{KL-Divergence Escalates Under Cognitive Load}

When forcing the network to operate under extreme truncation ($k=1$), we uncovered a severe dependency related to structural prompt complexity. 

The Kullback-Leibler (KL) divergence between baseline logit probabilities and $k$-truncated probabilities was calculated per taxonomy stage. 
As cognitive load transitions upward (i.e. moving from level 1 ``Remembering'' toward level 6 ``Creating''), the output logit distributions drift aggressively under tight rank ablation ($k \le 2$). Notably, divergence levels for simple factual recall tasks remain tightly bounded near zero even at $k=1$. In stark contrast, synthetic ``Evaluating'' logic forces significant KL divergence escalations under identical $k=1$ constraints. Standard divergence effectively vanishes across all taxonomic brackets uniformly when constraints relax to $k=4$, demonstrating an inherent limit to essential structural dimensions.

\subsection{Generative Match Rates (Fidelity)}

Divergence in theoretical logit spaces naturally translates to autoregressive destruction. We tracked exactly how often a truncated Phi-2 instance generated the verbatim identical string as the unaltered baseline model:

\begin{itemize}
    \item At $k=1$, the model sustained an $\sim10\%$ identical match rate for low-order \textit{Remembering} and \textit{Understanding} commands.
    \item At $k=1$, output fidelity collapsed to roughly $0\%$ when subjected to \textit{Creating} and \textit{Evaluating} BTL levels, resulting in hallucinatory cascades.
    \item At $k \ge 3$, the generative outputs matched standard outputs roughly $70-90\%$ of the time uniformly across the hierarchy.
\end{itemize}

These responses explicitly confirm that deep synthesis heavily depends on the minute variations afforded by trailing singular vectors ($v_3$ to $v_5$), whereas simple factual routing depends fundamentally only on the primary semantic driver ($v_1$).

\section{Inferences and Discussion}

The integration of these dynamics outlines several key characteristics of Large Language Model mechanics:

\paragraph{1. MHA is an Exploitable Low-Rank Dissipator}
The fundamental mathematical properties generated by MHA are heavily redundant. Because $k=3$ recovers around 95\% of spectral energy and preserves autoregressive rationality across all cognitive tasks, standard large-dimension multi-head attention over-calculates interactions. The bulk of token relations consolidate natively into low-rank manifolds.

\paragraph{2. The Purpose of Minor Singular Modes}
It is not merely that "more complex tasks" generate "more energy." Rather, tasks requiring abstraction or complex planning derive their structural coherence from long-tail combinations. Because basic semantic proximity is locked into the primary singular vector, $v_1$, eliminating subsequent minor modes destroys the subtle syntactical tracking that guides nuanced, multi-step Chain of Thought behaviors. 

\paragraph{3. Adaptive Inference Optimization Pathway}
Given that the KL-divergence effectively correlates with the hierarchical cognitive load of the prompt, model deployment infrastructures can exploit this variance. A gating mechanism could determine structural prompt complexity upstream; low-order information retrieval tasks could natively utilize structurally truncated operators, dramatically reducing memory bandwidth latency inside the matrix multiplication pathways, while dedicating full-rank resolution strictly to complex generative synthesis.

\section{Conclusion}

Our exhaustive SVD intervention sweep across Phi-2 verifies that the spectral properties of layer-wise Multi-Head Attention fundamentally shift in accordance directly with prompt-induced cognitive load. This interaction highlights that Transformer models mathematically scale their active representations. Standard attention channels function continuously as low-rank dissipating funnels, wherein basic semantic queries bypass structural limitations effectively, yet complex generative synthesis collapses when stripped of long-tail singular vectors. Ultimately, mapping taxonomy-level inference complexity to necessary effective rank reveals unprecedented opportunities to regulate inference speed organically in response to required theoretical rigor.

\end{document}
